\chapter{Electromagnetic Induction}

\section{Faraday's Law and Lenz's Law}

\textbf{Faraday's law} and \textbf{Lenz's law} are usually written together as
\begin{equation} \label{eq:faraday-lenz-law}
    \epsilon = - \odv{\Phi_B}{t}
\end{equation}
Strictly speaking, Faraday's law gives the magnitude,
while Lenz's law gives the sign.

\section{Maxwell's Equations}

\textbf{Maxwell's equations} are essentially the essence
of electricity and magnetism distilled into four equations.
You have more or less seen them already.
The equations are
\begin{subequations}
    \begin{equation}
        \oiint_S \vecb{E} \cdot \odif{\vecb{A}} = \frac{q}{\epsilon_0},
    \end{equation}
    \begin{equation}
        \oiint_S \vecb{B} \cdot \odif{\vecb{A}} = 0,
    \end{equation}
    \begin{equation}
        \oint_C \vecb{E} \cdot \odif{\vecb{\ell}} = - \odv{\Phi_B}{t},
    \end{equation}
    \begin{equation}
        \oint_C \vecb{B} \cdot \odif{\vecb{\ell}}
        = \mu_0 I + \mu_0 \epsilon \odv{\Phi_E}{t}.
    \end{equation}
\end{subequations}
Indeed, these are Gauss's law for electricity [\pgRef{eq:gausss-law}],
Gauss' law for magnetism, [\pgRef{eq:gausss-law-magnetism}]
the Faraday-Lenz law [\pgRef{eq:faraday-lenz-law}],
and Ampère's law [\pgRef{eq:amperes-law}] (sort of).
It turns out that Ampère's law was incomplete,
and Maxwell added the \(\mu_0 \epsilon \odv{\Phi_E}{t}\) term,
with the combined equation known as the \textbf{Ampère-Maxwell law}.

The Faraday-Lenz law suggests that changing magnetic fields
create electric fields,
and the Ampère-Maxwell law suggests that changing electricity fields
create magnetic fields.
This suggests that it is possible for these two effects to chain together
into oscillating electric and magnnetic fields,
known as \textbf{electromagnetic waves} (\textbf{EM waves}).
Maxwell correctly postulated that his mathematical electromagnetic waves
were in fact light.
From Maxwell's equations,
we can derive that the speed of light \(c\) is
\begin{equation}
    c = \frac{1}{\sqrt{\mu_0 \epsilon_0}}
\end{equation}
The exact derivation is too advanced for this volume
and will be provided in Vol.~3: Waves and Oscillations.
